\Article{15}
There remains, of course, the case in which the point $\mu$ is off the surface $t$ itself, yet at an imperceptible distance from it. In this case our integral can indeed have a perceptible value, which we shall now examine more closely.

Let the straight line drawn from the point $\mu$ normal to the surface $t$ and the fixed line issuing from the same point meet the spherical surface at the points $G$, $H$, respectively. Let the arc $GH=k$ and the arc between $G$ and an arbitrary point of the spherical surface be $=v$. Finally, let $w$ be the spherical angle between the arcs $k$, $v$. For the element $\mathrm{d}\Pi$ we may thus take the product $\sin v.\mathrm{d}v.\mathrm{d}w$. Writing $r$ for $(\mu,\mathrm{d}t)$ for brevity, the integral (II) becomes
\[
=\int\int\pm(\cos k.\cos v+\sin k.\sin v.\cos w)\theta r.\sin v.\mathrm{d}v.\mathrm{d}w
\]
This integration need extend only over those parts of the spherical surface onto which the imperceptible distances $r$ are projected. These correspond to the imperceptible part of the surface $t$. If we treat it as \textit{flat} and denote the least distance (corresponding to the point $G$ or the value $v=0$) by $\rho$, we get $r=\frac{\rho}{\cos v}$, which is independent of $w$. After integration with respect to $w$, say, from $w=0$ to $w=360^0$, the integral becomes
\[
=\pm\int2\pi\theta r.\cos k.\cos v.\sin v.\mathrm{d}v=\pm\int\frac{2\pi\cos k.\rho\rho\theta r.\mathrm{d}r}{r^3}
\]
This integration must extend from $r=\rho$ to an arbitrary perceptible value, however small. If, then, we set in general
\[
2rr\int\frac{\theta r.\mathrm{d}r}{r^3}=-\theta'r
\]
\SourceLeaf{55}choosing the integration constant so that $\int\frac{\theta r.\mathrm{d}r}{r^3}=0$ for an arbitrary perceptible value within the range for which we can perform experiments, the integral (II), with imperceptible quantities disregarded, will be
\[
=\pm\pi\cos k.\theta'\rho
\]

If it seems doubtful whether we may treat the part of the surface $t$ within an imperceptible distance of the point $\mu$ as flat, let us consider a spherical surface in its place. Let $R$ be the distance from the center of the sphere to the point $\mu$, taken as positive or negative according as the center is in the direction toward $G$ or in the opposite direction. Then
\[
\begin{gathered}
\cos v=\frac{\rho}{r}(1-\frac{\rho}{2R})+\frac{r}{2R}\\
\sin v.\mathrm{d}v=[\frac{\rho}{rr}(1-\frac{\rho}{2R})-\frac{1}{2R}]\mathrm{d}r
\end{gathered}
\]
It readily follows that the integral in this case does not differ by a perceptible amount from the value $\pm\pi\cos k.\theta'\rho$ found earlier, provided $R$ is a perceptible quantity. Now whatever the curvature of the surface $t$ in the part in question, provided the radii of curvature are perceptible, we can always specify two spherical surfaces tangent to the surface $t$ at its point nearest $\mu$, such that $t$ lies between them and their radii are of perceptible magnitude. Clearly our integral then lies between the integrals for those surfaces, and hence the same formula expresses it without perceptible error. The only exceptions occur when the surface $t$, within an imperceptible distance of the point $\mu$, presents either a curvature of imperceptible radius, an edge or a cusp.

\Article{16}
If we now proceed from integration (II) to integral (I), the latter clearly becomes imperceptible in either of two cases. The first is when the former has yielded a perceptible value at \textit{no} point of the surface $T$. The second is when the elements of the surface $T$ at whose points integral (II) became perceptible together form an area of only imperceptible magnitude. On due consideration it will be seen that integral (I) can acquire a perceptible value only insofar as the surface $T$ contains one or more parts of perceptible magnitude lying at an imperceptible distance from the surface $t$. Since such parts cannot depart perceptibly from being parallel to the surface $t$, at any point of them $\cos k$ will not differ percep\SourceLeaf{56}tibly from either $+1$ or $-1$, according as the outer or inner face of the surface $T$ is turned toward the surface $t$. Let $\NearTau$, $\NearTau'$ denote the parts of the surface $T$ at an imperceptible distance from the surface $t$, with $\NearTau$ denoting the parts where the outer face of one surface faces the inner face of the other, and $\NearTau'$ the parts where faces of the same kind face one another. Finally, let $\rho$ denote the least distance of any element $\mathrm{d}\NearTau$ or $\mathrm{d}\NearTau'$ from the surface $t$. Disregarding imperceptible quantities, our integral (I) becomes
\[
=-\int\pi\theta'\rho.\mathrm{d}\NearTau+\int\pi\theta'\rho.\mathrm{d}\NearTau'
\]
Clearly it makes no difference here whether the parts $\NearTau$, $\NearTau'$ are referred to the surface $T$ or to $t$.

We have thus obtained a complete solution of the problem we set ourselves in Article 6, for the property of the function $\varphi$ that underlies the main investigation into the equilibrium shape of fluids. Specifically, we have
\[
\int\int\mathrm{d}s.\mathrm{d}S.\varphi(\mathrm{d}s,\mathrm{d}S)=4\pi\sigma\psi0-\pi\tau\theta0+\pi\tau'\theta0-\pi\int\mathrm{d}\NearTau.\theta'\rho+\pi\int\mathrm{d}\NearTau'.\theta'\rho
\]

\Article{17}
The origin of the function $\theta'$ can be expressed by writing
\[
\frac{\theta'r}{rr}=\int\frac{2\theta x.\mathrm{d}x}{x^3}
\]
with the integral taken from $x=r$ up to an arbitrary perceptible constant value, which we denote here by $R$. This integral is clearly less than $\int\frac{2\theta r.\mathrm{d}x}{x^3}$ between the same limits, which is $=\frac{\theta r}{rr}-\frac{\theta r}{RR}$, and hence all the more so less than $\frac{\theta r}{rr}$. But since the indefinite integral is
\[
\int\frac{2\theta x.\mathrm{d}x}{x^3}=-\frac{\theta x}{xx}+\int\frac{\mathrm{d}\theta x}{xx}=-\frac{\theta x}{xx}-\int\frac{\psi x.\mathrm{d}x}{xx}
\]
we have
\[
\frac{\theta'r}{rr}=\frac{\theta r}{rr}-\frac{\theta R}{RR}-\int\frac{\psi x.\mathrm{d}x}{xx}
\]
with the integral taken between the same limits. It is less than the integral $\int\frac{\psi r\mathrm{d}x}{xx}$ and hence also less than $\frac{\psi r}{r}$. Therefore the value of $\frac{\theta'r}{rr}$ is greater than
\[
\frac{\theta r}{rr}-\frac{\theta R}{RR}-\frac{\psi r}{r}
\]
\SourceLeaf{57}Thus $\theta'r$ falls between the bounds
\[
\theta r\quad\text{and}\quad\theta r-rr.\frac{\theta R}{RR}-r\psi r
\]
whose difference can clearly become smaller than any assignable quantity as $r$ decreases indefinitely, since we suppose even $\psi0$ to be finite. From this we conclude that we must set $\theta'0=\theta0$. In the formula we arrived at in the preceding article, the term $-\pi\tau\theta0$ can evidently be regarded as included in the term $-\pi\int\mathrm{d}\NearTau.\theta'\rho$, and the term $\pi\tau'\theta0$ as included in the term $\pi\int\mathrm{d}\NearTau'.\theta'\rho$, if we remove the distinction we have made between an imperceptible distance and zero distance and count the parts $\tau$, $\tau'$ with the parts $\NearTau$, $\NearTau'$, respectively. Yet although this would make the solution more elegant in a mathematical sense, for our purpose it is preferable to retain the distinction.

